% Shared macros.
\newcommand{\R}{\mathbb R}
\newcommand{\N}{\mathbb N}
\newcommand{\E}{\mathbb E}
\newcommand{\PP}{\mathbb P}
\newcommand{\dd}{\,\mathrm d}
\newcommand{\eps}{\varepsilon}
\newcommand{\norm}[1]{\lVert #1\rVert}
\newcommand{\abs}[1]{\lvert #1\rvert}
\newcommand{\TV}{\mathrm{TV}}
\DeclareMathOperator{\Var}{Var}
\DeclareMathOperator{\Cov}{Cov}
\DeclareMathOperator{\tr}{tr}
\DeclareMathOperator{\supp}{supp}
\DeclareMathOperator{\Ent}{Ent}
\DeclareMathOperator{\Law}{Law}

% Label aliases: \aliaslabel{new}{old} makes \ref/\cref{new} resolve to the label old
% (used where a source label was kept on a statement that already carries another label).
\makeatletter
\newcommand{\aliaslabel}[2]{\AtBeginDocument{%
  \@ifundefined{r@#2}{}{%
    \expandafter\let\csname r@#1\expandafter\endcsname\csname r@#2\endcsname
    \expandafter\let\csname r@#1@cref\expandafter\endcsname\csname r@#2@cref\endcsname}}}
\makeatother
\aliaslabel{b:eq:SK-covariance}{eq:common-covariance}
\aliaslabel{f:completion}{eq:model-completion-laplace}
\aliaslabel{b:eq:os-contacts}{b:eq:contacts}

% Macros for Part IV (see NOTATION.md).
\DeclareMathOperator{\KL}{KL}
\newcommand{\ipn}[2]{\langle #1,#2\rangle}
\newcommand{\normn}[1]{\lVert #1\rVert}
\newcommand{\ex}[1]{\mathsf{#1}}
\DeclareMathOperator{\Up}{Up}
\newcommand{\Upt}{\Up_{1/2}}
\DeclareMathOperator{\Unif}{Unif}
